\conclusions\label{ch:concl}
    \par
        Questa Tesi di laurea mirava a fare luce sulla figura e sul ruolo del Tecnico Informatico all'interno di un Polo Universitario, ruolo che spesso viene trascurato nel settore nonostante sia un ruolo chiave per garantire la continuità dei servizi TIC offerti dalle Università, oltre che da qualsiasi altra istituzione, in un mondo ormai sempre più digitalizzato e tecnologico. 
    \par
        Durante questo periodo di tirocinio ho vissuto a stretto contatto con la Direzione del Centro Polifunzionale di Gorizia, interfacciandomi anche con i tecnici del DAGS per poter fare interventi specifici su loro richiesta. Questa esperienza mi ha permesso di lavorare in un ambiente familiare ma, da un altro punto di vista, mi ha permesso anche di approfondire scenari e situazioni che prima mi erano sconosciuti. Il tirocinio mi ha permesso di mettere in gioco le mie abilità non solo tecniche ma anche di adattamento e di gestione degli interventi, in quanto molto spesso le decisioni tecniche dovevano essere prese nell'arco di qualche minuto. \newline
        Le mansioni ricoperte erano state precedentemente illustrate in occasione di un colloquio conoscitivo avvenuto prima dell'avvio del tirocinio, tuttavia queste si sono ampliate in corso d'opera. Il Corso di Studi in Scienze e Tecnologie Multimediali, attualmente Corso di Studi in Tecnologie digitali e comunicazione per le industrie creative, mi ha fornito delle conoscenze più che approfondite in molte aree e mi ha permesso di affrontare tutti i compiti richiesti e ciò che di nuovo ho potuto sperimentare. 
    \par
        Dato il rapido sviluppo delle nuove tecnologie, nello specifico dell'Intelligenza Artificiale, mi è sembrato opportuno discutere della figura del Tecnico anche ponendola a confronto con le IA e di come questi due diversi soggetti possano convivere ed esistere. Le IA offrono una maggiore rapidità di risposta a discapito però dell'accuratezza, mentre il Tecnico Informatico può intervenire sempre, garantendo il corretto funzionamento del dispositivo. 
        Da questo tirocinio ho appreso inoltre che la componente fondamentale di qualsiasi intervento, oltre alle conoscenze pregresse, è il sapersi interfacciare con l'ambiente umano e l'empatia che ne deriva; due fattori chiave che nonostante la IA possa provare a replicare, non saranno mai al pari di un qualsiasi rapporto e intervento con un Tecnico vero e proprio. 
    \par
        Infine questa esperienza mi ha permesso di accrescere le mie abilità di Tecnico Informatico, specializzandomi in interventi in loco in ambito informatico-multimediale. Il tirocinio inoltre mi ha aiutato a fare chiarezza sul mio futuro, orientandomi in direzione del lavoro preferibilmente presso le Pubbliche Amministrazioni, senza tuttavia escludere a priori l'ambito privato.